% Generated by scripts/gen_blueprint.py from blueprint/nodes/06-four-exponentials.toml. Do not edit.

\chapter{The four exponentials theorem in transcendence degree one}
\label{chap:06-four-exponentials}

The four exponentials conjecture holds when the four logarithms generate a field of transcendence degree one: Theorem~\ref{thm:four_exponentials_trdeg_one}, proved independently by Waldschmidt and by Brownawell. The proof here is Waldschmidt's of 1973 \cite{W73} with the tools of his 1971 paper \cite{W71}, so it needs neither a lower bound for linear forms in logarithms nor Philippon's zero estimate. The printed 1973 proof meets one hypothesis of its Lemme~2 with Gel'fond's bound for linear forms in two logarithms; the zero count of Theorem~\ref{thm:expPoly_zero_count}, which needs no separation between the frequencies, makes that hypothesis unnecessary.

The chapter has three parts: Gel'fond's transcendence criterion in Waldschmidt's form, the zero count, and the construction of 1973. The construction proves by contradiction that its hypotheses cannot hold together, so several statements of the last part have hypotheses that are never satisfied.

\textbf{Notation for the construction.} $N$ is a large integer, and
\[ S=\Bigl\lfloor\frac{N^{2}}{\sqrt{\log N}}\Bigr\rfloor,\qquad S'=\lfloor S/2\rfloor,\qquad t_1=\Bigl\lfloor\frac{N}{\sqrt{\log N}}\Bigr\rfloor,\qquad t_2=\bigl\lfloor N\sqrt{\log N}\bigr\rfloor, \]
\[ R_1=14t_1,\qquad R_2=14t_2. \]
For complex coefficients $c(i,j,h)$, $i<S$ and $j,h<2N$, the auxiliary function is
\[ F(z)=\sum_{i<S}\ \sum_{j,h<2N}c(i,j,h)\,z^{i}e^{(jx_1+hx_2)z}. \]
\emph{Presentation data} for $x_1,x_2,y_1,y_2\in\C$ are: $\omega,\omega_1\in\C$ with $\omega$ transcendental; $Q\in\Z[X][Y]$, monic in $Y$ of degree $d\ge1$, with $Q(\omega,\omega_1)=0$ and minimal there (no non-zero $A\in\Z[X][Y]$ with $\deg_YA<d$ vanishes at $(\omega,\omega_1)$); and $D,E_i,G_j,H_{ij}\in\Z[X][Y]$ with $D(\omega,\omega_1)\neq0$ and
\[ x_iD=E_i,\qquad y_jD=G_j,\qquad e^{x_iy_j}D=H_{ij}\qquad\text{at }(\omega,\omega_1). \]
For $M\in\N$ and integers $q(i,j,h,\mu,\nu)$, $\mu<M$, $\nu<d$, put $c_q(i,j,h)=\sum_{\mu<M,\,\nu<d}q(i,j,h,\mu,\nu)\,\omega^{\mu}\omega_1^{\nu}$, and let $F_q$ be the auxiliary function with these coefficients (with $T$ in place of $2N$ where a statement says so). The growth functions with constant $k>0$ are
\[ \sigma_1(x)=k\cdot\begin{cases}x-3+9\sqrt{\log3},&x\le3,\\ x^{2}\sqrt{\log x},&x>3,\end{cases}\qquad\sigma_2(x)=k\cdot\begin{cases}x-3+9/\sqrt{\log3},&x\le3,\\ x^{2}/\sqrt{\log x},&x>3.\end{cases} \]
Given functions $\sigma_1,\sigma_2$, a polynomial $P\in\Z[X]$ is \emph{small at $\omega$ for $(N,C)$} if $P\neq0$, all its coefficients have absolute value at most $e^{\sigma_1(N)}$, $\deg P\le\sigma_2(N)$, and $|P(\omega)|<e^{-C\sigma_1(N)\sigma_2(N)}$.

\section{Gel'fond's transcendence criterion}

% Prove2Me: FourExp.height_dvd_le -- https://prove2.me/theorems/5b99493b-fb4c-49c3-8afe-4ed79a8f2bed
\begin{lemma}[Height of a divisor]
\label{lem:height_dvd_le}
\lean{Diaz.height_dvd_le}
\leanok
Let $P\in\Z[X]$ be non-zero and let $Q\in\Z[X]$ divide $P$. If every coefficient of $P$ has absolute value at most $H$, then every coefficient of $Q$ has absolute value at most $e^{\deg P}H$.

{\small\emph{Source:} \cite[\S3, Lemme 1]{W71} (Gel'fond; a generalisation of a lemma of Popken and Koksma).}
\end{lemma}

\begin{proof}
\leanok
Gel'fond's inequality $H(P_1)H(P_2)\le e^{\deg(P_1P_2)}H(P_1P_2)$, where $H$ is the largest absolute value of a coefficient, follows by comparing heights with Mahler measure. Apply it to $P=QR$, where $H(R)\ge1$.
\end{proof}

% Prove2Me: FourExp.small_irreducible_factor -- https://prove2.me/theorems/0d10482b-a0ff-46c6-a565-57809008dcf0
\begin{lemma}[Gel'fond's lemma]
\label{lem:small_irreducible_factor}
\lean{Diaz.small_irreducible_factor}
\leanok
Let $\alpha\in\C$ be transcendental and let $P\in\Z[X]$ be primitive, with every coefficient of absolute value at most $H$. Let $n,\lambda\in\R$ with $\deg P\le n\le\log H$ and $\lambda>6$. If $|P(\alpha)|<H^{-\lambda n}$, then there are a primitive irreducible divisor $Q$ of $P$ and an integer $s\ge1$ with
\[ |Q(\alpha)|<H^{-(\lambda-6)n/s},\qquad\deg Q\le n/s, \]
and every coefficient of $Q$ of absolute value at most $H^{1/s}e^{2n/s}$.

{\small\emph{Source:} \cite[\S3, Lemme 2]{W71}, after Gel'fond \cite{Gel52}.}
\end{lemma}

\begin{proof}
\leanok
\uses{lem:exists_irreducible_factor_cofactor_bound}
Write $P=Q^{s}R$ as in Lemma~\ref{lem:exists_irreducible_factor_cofactor_bound}. The cofactor $R$ is not too small at $\alpha$, so $|Q(\alpha)|^{s}$ is small; the height and degree of $Q$ follow through Mahler measure.
\end{proof}

% Prove2Me: FourExp.dvd_of_small_values -- https://prove2.me/theorems/94ad0017-5959-42e3-8dbb-5279735f4026
\begin{lemma}[Small values force divisibility]
\label{lem:dvd_of_small_values}
\lean{Diaz.dvd_of_small_values}
\leanok
Let $P,Q\in\Z[X]$ with $Q$ irreducible, let $\alpha\in\C$, and let $H,h\ge1$ bound the absolute values of the coefficients of $P$ and of $Q$ respectively. Put $d=\deg P$ and $\delta=\deg Q$. If
\[ \bigl((1+|\alpha|)(d+\delta)\bigr)^{d+\delta}H^{\delta}h^{d}\bigl(|P(\alpha)|+|Q(\alpha)|\bigr)<1, \]
then $Q$ divides $P$.

{\small\emph{Source:} \cite[\S3, (3.13)]{W71}, citing \cite[Ch.~V, \S2]{Lang66}.}
\end{lemma}

\begin{proof}
\leanok
Otherwise $P$ and $Q$ are coprime, and their resultant is a non-zero integer of the form $AP+BQ$ with $A,B\in\Z[X]$ bounded by Hadamard's inequality; evaluating at $\alpha$ contradicts the hypothesis.
\end{proof}

% Prove2Me: FourExp.dvd_of_small_values_at_scale -- https://prove2.me/theorems/d720908d-8850-4491-a909-d3cc6d5a732a
\begin{lemma}[Divisibility at one scale]
\label{lem:dvd_of_small_values_at_scale}
\lean{Diaz.dvd_of_small_values_at_scale}
\leanok
Let $\alpha\in\C$ and $\varepsilon>0$. There is $U$ such that for all real $u\ge U$ and $1\le v\le u$ the following holds. Let $P,Q\in\Z[X]$ with $Q$ irreducible, the coefficients of $P$ of absolute value at most $e^{u}$ and $\deg P\le v$, the coefficients of $Q$ of absolute value at most $e^{3u}$ and $\deg Q\le(1+\varepsilon/2)v$, and $|P(\alpha)|<e^{-(4+\varepsilon)uv}$, $|Q(\alpha)|<e^{-(4+\varepsilon)uv}$. Then $Q$ divides $P$.

{\small\emph{Source:} A step of the proof of \cite[\S3, Lemme fondamental]{W71}.}
\end{lemma}

\begin{proof}
\leanok
\uses{lem:dvd_of_small_values}
Lemma~\ref{lem:dvd_of_small_values}, with every size written in terms of $(u,v)$.
\end{proof}

% Prove2Me: FourExp.transcendence_criterion_continuous -- https://prove2.me/theorems/edd2f289-6d1c-476d-ad9a-5b7015839de3
\begin{lemma}[The criterion for continuous growth functions]
\label{lem:transcendence_criterion_continuous}
\lean{Diaz.transcendence_criterion_continuous}
\leanok
Let $\alpha\in\C$ and $\varepsilon>0$. Let $\sigma_1,\sigma_2\colon\R\to\R$ be continuous, strictly increasing and tending to $+\infty$, and let $a_1,a_2\ge1$ be such that for every $x\ge1$
\[ \sigma_2(x)\le\sigma_1(x),\qquad\sigma_i(x+1)\le a_i\,\sigma_i(x)\quad(i=1,2). \]
Put $C=\max\{10+\varepsilon,(4+\varepsilon)a_1a_2\}$. Suppose that for every integer $N>N_0$ there is a non-zero $P_N\in\Z[X]$ with every coefficient of absolute value at most $e^{\sigma_1(N)}$, $\deg P_N\le\sigma_2(N)$ and $|P_N(\alpha)|<e^{-C\sigma_1(N)\sigma_2(N)}$. Then $\alpha$ is algebraic.

{\small\emph{Source:} \cite[\S3, Lemme fondamental]{W71} and its proof, where continuity of the $\sigma_i$ is assumed without loss of generality.}
\end{lemma}

\begin{proof}
\leanok
\uses{lem:height_dvd_le,lem:small_irreducible_factor,lem:dvd_of_small_values_at_scale}
Waldschmidt's proof. If $\alpha$ is transcendental, Lemma~\ref{lem:small_irreducible_factor} gives small irreducible factors of the $P_N$. At a suitable scale Lemma~\ref{lem:dvd_of_small_values_at_scale} makes such a factor divide a later $P_N$, and Lemma~\ref{lem:height_dvd_le} then bounds its height, against the choice of the scale.
\end{proof}

% Prove2Me: FourExp.transcendence_criterion -- https://prove2.me/theorems/65f053a5-33bf-46c5-95a8-ac81974255c5
\begin{theorem}[Waldschmidt's transcendence criterion]
\label{thm:transcendence_criterion}
\lean{Diaz.transcendence_criterion}
\leanok
Let $\alpha\in\C$ and $\varepsilon>0$. Let $\sigma_1,\sigma_2\colon\R\to\R$ be strictly increasing and tending to $+\infty$, and let $a_1,a_2\ge1$ be such that for every $x>0$
\[ \sigma_2(x)\le\sigma_1(x),\qquad\sigma_i(x+1)\le a_i\,\sigma_i(x)\quad(i=1,2). \]
Put $C=\max\{10+\varepsilon,(4+\varepsilon)a_1a_2\}$. Suppose that for every integer $N>N_0$ there is a non-zero $P_N\in\Z[X]$ with every coefficient of absolute value at most $e^{\sigma_1(N)}$, $\deg P_N\le\sigma_2(N)$ and $|P_N(\alpha)|<e^{-C\sigma_1(N)\sigma_2(N)}$. Then $\alpha$ is algebraic.

{\small\emph{Source:} \cite[\S3, Lemme fondamental]{W71}, with constant $a_i$ as in Remark 2 there; a refinement of Gel'fond's criterion \cite[Ch.~III, \S4, Lemma VII]{Gel52}.}
\end{theorem}

\begin{proof}
\leanok
\uses{lem:transcendence_criterion_continuous}
Replace $\sigma_1,\sigma_2$ by their piecewise linear interpolations between the integers and apply Lemma~\ref{lem:transcendence_criterion_continuous}.
\end{proof}

\section{Zeros of exponential polynomials}

% Prove2Me: FourExp.expPoly_value_le_derivs -- https://prove2.me/theorems/aec4da14-b097-4fda-999a-687c2dbecdd4
\begin{lemma}[Values from low-order derivatives]
\label{lem:expPoly_value_le_derivs}
\lean{Diaz.expPoly_value_le_derivs}
\leanok
Let $w_1,\dots,w_l\in\C$ be pairwise distinct with $|w_j|\le W$ for some $W\ge0$. For each $j$ let $P_j\in\C[z]$ be zero or of degree less than $q_j$, and put $g(z)=\sum_jP_j(z)e^{w_jz}$ and $n=\sum_jq_j$. Let $R\ge0$. If $|g^{(s)}(0)|\le D$ for every $s<n$, then
\[ |g(u)|\le n\,(W+1)^{n+1}e^{R(W+1)}D\qquad\text{whenever }|u|\le R. \]

{\small\emph{Source:} \cite[\S4, (4.5)--(4.13)]{W71}.}
\end{lemma}

\begin{proof}
\leanok
\uses{lem:expPoly_iteratedDeriv_le}
Expand $g$ in its Taylor series at $0$ and bound every derivative there with Lemma~\ref{lem:expPoly_iteratedDeriv_le}.
\end{proof}

% Prove2Me: FourExp.expPoly_zero_count_scaled -- https://prove2.me/theorems/0fdacfda-2630-482e-8cf1-f159d6ba6be7
\begin{lemma}[The rescaled inequality]
\label{lem:expPoly_zero_count_scaled}
\lean{Diaz.expPoly_zero_count_scaled}
\leanok
Let $f(z)=\sum_{j=1}^{l}\sum_{i<q_j}b_{j,i}z^{i}e^{\omega_jz}$ with pairwise distinct $\omega_j$ and $b_{j,i}$ not all zero; put $n=\sum_jq_j$ and $\Omega=\max_j|\omega_j|$, and suppose $\Omega>0$. Let $z_0\in\C$, $\rho\ge0$, let $S$ be a finite set of points of the disc $|z-z_0|\le\rho$, and put $\sigma=\sum_{z\in S}\ord_zf$ and $x=\rho\Omega$. Then for every $R>x+1$
\[ \sigma\log\frac{R-x}{x+1}\le\log\Bigl(n!\,2^{n+1}\frac{R}{R-1}\Bigr)+2R. \]

{\small\emph{Source:} \cite[\S4, (4.14)]{W71} and the rescaling that follows it.}
\end{lemma}

\begin{proof}
\leanok
\uses{lem:cauchy_estimate_with_zeros,lem:expPoly_value_le_derivs,lem:expPoly_ne_zero}
Lemmas~\ref{lem:cauchy_estimate_with_zeros} and~\ref{lem:expPoly_value_le_derivs} bound $\max_{s<n}|g^{(s)}(0)|$ against the maximum of $|g|$ on $|u|=R$ in opposite directions; apply both to $g(z)=f(z_0+z/\Omega)$, which is not identically zero by Lemma~\ref{lem:expPoly_ne_zero}.
\end{proof}

% Prove2Me: FourExp.zero_count_arith -- https://prove2.me/theorems/52059e8b-dd9e-404c-ab6e-cd03a25e3a57
\begin{lemma}[From the rescaled inequality to the count]
\label{lem:zero_count_arith}
\lean{Diaz.zero_count_arith}
\leanok
Let $n\ge2$ be an integer and let $x\ge0$, $\lambda>0$ and $\sigma\ge0$ be real. If for every $R>x+1$
\[ \sigma\log\frac{R-x}{x+1}\le\log\Bigl(n!\,2^{n+1}\frac{R}{R-1}\Bigr)+2R, \]
then
\[ \sigma<\frac{n}{\lambda}+2\,\frac{1+n^{\lambda}}{\lambda\log n}\,(1+x). \]

{\small\emph{Source:} After \cite[\S4, (4.14)]{W71}.}
\end{lemma}

\begin{proof}
\leanok
For most parameters take $R=n^{\lambda}(1+x)+x$, so that $\log\frac{R-x}{x+1}=\lambda\log n$. The remaining cases, among them $2\le n\le5$, where the bound $n!\,2^{n}\le n^{n}$ used in \cite{W71} fails, need other choices of $R$.
\end{proof}

% Prove2Me: FourExp.zero_count_arith_poly -- https://prove2.me/theorems/46d2beb4-7606-4504-994b-05af059e4f49
\begin{lemma}[Few zeros beat the bound]
\label{lem:zero_count_arith_poly}
\lean{Diaz.zero_count_arith_poly}
\leanok
For every integer $n\ge1$, real $x\ge0$ and $\lambda>0$, and every integer $0\le\sigma\le n-1$,
\[ \sigma<\frac{n}{\lambda}+2\,\frac{1+n^{\lambda}}{\lambda\log n}\,(1+x), \]
where for $n=1$ the second term is read as $0$.

{\small\emph{Source:} Elementary.}
\end{lemma}

\begin{proof}
\leanok
For $n=1$, $\sigma=0<1/\lambda$; for $\lambda\le1$, $n/\lambda\ge n>\sigma$. For $\lambda>1$, study the right-hand side as a function of $\lambda\log n$.
\end{proof}

% Prove2Me: FourExp.zero_count_degenerate -- https://prove2.me/theorems/e7bb8de6-eb16-421c-a71b-a457db838e41
\begin{lemma}[The degenerate cases]
\label{lem:zero_count_degenerate}
\lean{Diaz.zero_count_degenerate}
\leanok
Let $f$, $n$ and $\Omega$ be as in Lemma~\ref{lem:expPoly_zero_count_scaled}, without the hypothesis $\Omega>0$, and suppose that $n\le1$ or $\Omega=0$. Then for every finite set $S\subset\C$
\[ \sum_{z\in S}\ord_zf\le n-1. \]

{\small\emph{Source:} Elementary; the degenerate cases of \cite[\S4, Lemme 3]{W71}.}
\end{lemma}

\begin{proof}
\leanok
If $n=1$, then $f=be^{\omega z}$ with $b\neq0$, which has no zeros. If $\Omega=0$, the only frequency is $0$ and $f$ is a non-zero polynomial of degree less than $n$.
\end{proof}

% Prove2Me: FourExp.expPoly_zero_count -- https://prove2.me/theorems/4d055466-a546-41a8-a433-78e09b040998
\begin{theorem}[Zeros of an exponential polynomial in a disc]
\label{thm:expPoly_zero_count}
\lean{Diaz.expPoly_zero_count}
\leanok
Let $\omega_1,\dots,\omega_l\in\C$ be pairwise distinct, let $q_1,\dots,q_l\in\N$, and let $b_{j,i}\in\C$ ($1\le j\le l$, $0\le i<q_j$) be not all zero. Put
\[ f(z)=\sum_{j=1}^{l}\sum_{i=0}^{q_j-1}b_{j,i}\,z^{i}e^{\omega_jz},\qquad n=\sum_jq_j,\qquad\Omega=\max_j|\omega_j|. \]
Let $z_0\in\C$, $\rho\ge0$ and $\lambda>0$. Then for every finite set $S$ of points of the disc $|z-z_0|\le\rho$
\[ \sum_{z\in S}\ord_zf<\frac{n}{\lambda}+2\,\frac{1+n^{\lambda}}{\lambda\log n}\,(1+\rho\Omega), \]
where for $n=1$ the second term is read as $0$. No separation between the $\omega_j$ is assumed.

{\small\emph{Source:} \cite[\S4, Lemme 3, (4.2)]{W71}; compare \cite{Tij71}.}
\end{theorem}

\begin{proof}
\leanok
\uses{lem:expPoly_zero_count_scaled,lem:zero_count_degenerate,lem:zero_count_arith,lem:zero_count_arith_poly}
For $n\ge2$ and $\Omega>0$, Lemmas~\ref{lem:expPoly_zero_count_scaled} and~\ref{lem:zero_count_arith}; otherwise Lemmas~\ref{lem:zero_count_degenerate} and~\ref{lem:zero_count_arith_poly}.
\end{proof}

% Prove2Me: FourExp.nonvanishing_derivative -- https://prove2.me/theorems/6b3e7ded-6c39-4625-b4ac-f87189de2224
\begin{lemma}[A non-zero derivative on the grid]
\label{lem:nonvanishing_derivative}
\lean{Diaz.nonvanishing_derivative}
\leanok
Let $x_1,x_2$ be linearly independent over $\Q$, and likewise $y_1,y_2$. Let $S,T,R_1,R_2,S'\in\N$ and let $c(i,j,h)\in\C$ ($i<S$, $j,h<T$) be not all zero; put
\[ G(z)=\sum_{i<S}\sum_{j,h<T}c(i,j,h)\,z^{i}e^{(jx_1+hx_2)z},\qquad n=ST^{2}. \]
If for some $\lambda>0$
\[ \frac{n}{\lambda}+2\,\frac{1+n^{\lambda}}{\lambda\log n}\Bigl(1+(R_1|y_1|+R_2|y_2|)\,T(|x_1|+|x_2|)\Bigr)\le R_1R_2S', \]
then $G^{(s)}(ay_1+by_2)\neq0$ for some $a<R_1$, $b<R_2$ and $s<S'$.

{\small\emph{Source:} \cite[Lemme 6]{W73}, with the zero estimate \cite[\S4, Lemme 3]{W71} in place of Gel'fond's zero lemma.}
\end{lemma}

\begin{proof}
\leanok
\uses{thm:expPoly_zero_count,lem:expPoly_ne_zero}
Otherwise the $R_1R_2$ points $ay_1+by_2$, distinct because $y_1,y_2$ are independent, would be zeros of order at least $S'$ of $G$. The frequencies $jx_1+hx_2$ are distinct, so $G$ is not identically zero (Lemma~\ref{lem:expPoly_ne_zero}), and Theorem~\ref{thm:expPoly_zero_count} counts fewer zeros.
\end{proof}

\section{The construction of 1973}

% Prove2Me: FourExp.rank_one_parametrization -- https://prove2.me/theorems/30f831f0-d392-4e30-aa69-dee7d4b3ba65
\begin{lemma}[A rank-one matrix of logarithms]
\label{lem:rank_one_parametrization}
\lean{Diaz.rank_one_parametrization}
\leanok
Let $l_{11},l_{12},l_{21},l_{22}$ be non-zero complex numbers with every $e^{l_{ij}}$ algebraic, $l_{11}l_{22}=l_{12}l_{21}$ and $\trdeg_{\Q}\Q[l_{11},l_{12},l_{21},l_{22}]\le1$, and suppose that neither the rows nor the columns of $(l_{ij})$ are linearly dependent over $\Q$. Then there are $x_1,x_2,y_1,y_2\in\C$, with $x_1,x_2$ linearly independent over $\Q$ and likewise $y_1,y_2$, such that every $e^{x_iy_j}$ is algebraic and $\trdeg_{\Q}\Q[x_1,x_2,y_1,y_2]\le1$.

{\small\emph{Source:} Elementary; as in \cite[\S I]{W73}.}
\end{lemma}

\begin{proof}
\leanok
Take $x=(l_{11},l_{21})$ and $y=(1,l_{12}/l_{11})$, so that $x_iy_j=l_{ij}$.
\end{proof}

% Prove2Me: FourExp.trdeg_one_presentation -- https://prove2.me/theorems/f300d712-2c41-46f4-99a9-5475f13944ac
\begin{lemma}[Presenting the field]
\label{lem:trdeg_one_presentation}
\lean{Diaz.trdeg_one_presentation}
\leanok
Let $x_1,x_2$ be linearly independent over $\Q$, and likewise $y_1,y_2$. If every $e^{x_iy_j}$ is algebraic and $\trdeg_{\Q}\Q[x_1,x_2,y_1,y_2]\le1$, then $x_1,x_2,y_1,y_2$ have presentation data.

{\small\emph{Source:} The opening of \cite[\S III]{W73}.}
\end{lemma}

\begin{proof}
\leanok
\uses{thm:hermite_lindemann_holds,lem:exists_monic_integral_model_presentation}
The number $\omega=x_1y_1$ is non-zero with $e^{\omega}$ algebraic, so it is transcendental by Theorem~\ref{thm:hermite_lindemann_holds}. Every $x_i$, $y_j$ and $e^{x_iy_j}$ is then algebraic over $\Q(\omega)$, and Lemma~\ref{lem:exists_monic_integral_model_presentation} presents them.
\end{proof}

% Prove2Me: FourExp.exists_pow_mul_pow_le_exp_sq_mul_sqrt_log -- https://prove2.me/theorems/1cc8c606-255e-418d-a988-71ed7b28748d
\begin{lemma}[The size estimate]
\label{lem:exists_pow_mul_pow_le_exp_sq_mul_sqrt_log}
\lean{Diaz.exists_pow_mul_pow_le_exp_sq_mul_sqrt_log}
\leanok
For all $c,r\in\N$ there is $\kappa>0$ such that for all integers $N\ge3$ and $S,T,a,b,m\in\N$ with
\[ S\le\frac{N^{2}}{\sqrt{\log N}},\quad m\le S,\quad T\le rN,\quad a\le r\frac{N}{\sqrt{\log N}},\quad b\le rN\sqrt{\log N}, \]
one has
\[ c^{\,c(1+m+S+T(a+b))}(1+m+S+T+a+b)^{c(1+m+S)}\le e^{\kappa N^{2}\sqrt{\log N}}. \]

{\small\emph{Source:} The parameter estimates in the proofs of \cite[Lemmes 4 and 7]{W73}.}
\end{lemma}

\begin{proof}
\leanok
\end{proof}

% Prove2Me: FourExp.exists_iteratedDeriv_presentation -- https://prove2.me/theorems/014c4a25-8d7b-43c6-b2ea-cdf3dab3ca82
\begin{lemma}[Derivatives at a grid point, presented]
\label{lem:exists_iteratedDeriv_presentation}
\lean{Diaz.exists_iteratedDeriv_presentation}
\leanok
Let $x_1,x_2,y_1,y_2\in\C$ with every $e^{x_iy_j}$ algebraic. Let $\omega,\omega_1\in\C$, let $\varphi$ be evaluation at $(\omega,\omega_1)$, and let $D,E_i,G_j,H_{ij}\in\Z[X][Y]$ with $\varphi(D)\neq0$, $x_i\varphi(D)=\varphi(E_i)$, $y_j\varphi(D)=\varphi(G_j)$ and $e^{x_iy_j}\varphi(D)=\varphi(H_{ij})$. There is $c\in\N$ such that for all $S,T,a,b,m\in\N$ there are $\Lambda\in\C$, $\Lambda\neq0$, with $|\Lambda|\le c^{\,c(1+m+S+T(a+b))}$, and $P_{ijh}\in\Z[X][Y]$ ($i<S$, $j,h<T$) of degree at most $c(1+m+S)$ in $X$ and in $Y$ and of length at most
\[ c^{\,c(1+m+S+T(a+b))}\,(1+m+S+T+a+b)^{c(1+m+S)}, \]
such that for all $f_{ijh}\in\C$
\[ \Lambda\cdot\frac{d^{m}}{dz^{m}}\Bigl(\sum_{i,j,h}f_{ijh}\,z^{i}e^{(jx_1+hx_2)z}\Bigr)\Big|_{z=ay_1+by_2}=\sum_{i,j,h}f_{ijh}\,\varphi(P_{ijh}). \]

{\small\emph{Source:} A step of the proofs of \cite[Lemmes 4 and 7]{W73}.}
\end{lemma}

\begin{proof}
\leanok
\uses{lem:exists_int_pow_repr,lem:length_mul_le,lem:length_sum_le}
Expand the derivative by Leibniz's rule. The powers of the algebraic numbers $e^{x_iy_j}$ reduce through integer relations (Lemma~\ref{lem:exists_int_pow_repr}), and lengths are tracked with Lemmas~\ref{lem:length_sum_le} and~\ref{lem:length_mul_le}.
\end{proof}

% Prove2Me: FourExp.exists_iteratedDeriv_reduced_presentation -- https://prove2.me/theorems/e8fbf45e-5c3e-460b-8f33-429a3e02cc84
\begin{lemma}[Reduced presentations]
\label{lem:exists_iteratedDeriv_reduced_presentation}
\lean{Diaz.exists_iteratedDeriv_reduced_presentation}
\leanok
In the setting of Lemma~\ref{lem:exists_iteratedDeriv_presentation}, let moreover $Q\in\Z[X][Y]$ be monic in $Y$ with $\varphi(Q)=0$. There is $c\in\N$ such that for all $S,T,M,a,b,m\in\N$ there are $\Lambda\neq0$ with $|\Lambda|\le c^{\,c(1+m+S+T(a+b))}$ and $R_{ijh\mu\nu}\in\Z[X][Y]$ ($i<S$, $j,h<T$, $\mu<M$, $\nu<\deg_YQ$) of $Y$-degree less than $\deg_YQ$ and $X$-degree at most $M+c(1+m+S)$, and of length at most
\[ c^{\,c(1+m+S+T(a+b))}\,(1+m+S+T+a+b)^{c(1+m+S)}, \]
such that for every integer array $q$
\[ \Lambda\,F_q^{(m)}(ay_1+by_2)=\varphi\Bigl(\sum_{i,j,h,\mu,\nu}q(i,j,h,\mu,\nu)\,R_{ijh\mu\nu}\Bigr), \]
where $F_q$ is formed with $T$ in place of $2N$.

{\small\emph{Source:} A step of the proofs of \cite[Lemmes 4 and 7]{W73}.}
\end{lemma}

\begin{proof}
\leanok
\uses{lem:exists_iteratedDeriv_presentation,lem:exists_modByMonic_length_le,lem:length_mul_le}
Put the coefficients $c_q$ into Lemma~\ref{lem:exists_iteratedDeriv_presentation} and reduce modulo $Q$ with Lemma~\ref{lem:exists_modByMonic_length_le}.
\end{proof}

% Prove2Me: FourExp.aux_linear_system -- https://prove2.me/theorems/278448f6-4953-4e99-acc6-da2040046df5
\begin{lemma}[The linear system]
\label{lem:aux_linear_system}
\lean{Diaz.aux_linear_system}
\leanok
Let $x_1,x_2,y_1,y_2\in\C$ have presentation data, and suppose every $e^{x_iy_j}$ is algebraic. There is $\kappa_1>0$ such that for every large $N$ there are an integer $M$ with $0<M\le\kappa_1S$ and an $R\times U$ integer matrix $B$, where $U=S(2N)^{2}Md$ and $2R\le U$, with entries of absolute value at most $e^{\kappa_1N^{2}\sqrt{\log N}}$, such that for every integer vector $q$ with $Bq=0$
\[ F_q^{(m)}(ay_1+by_2)=0\qquad\text{for all }a<t_1,\ b<t_2,\ m<S. \]

{\small\emph{Source:} The equation count of \cite[Lemme 4]{W73}.}
\end{lemma}

\begin{proof}
\leanok
\uses{lem:exists_iteratedDeriv_reduced_presentation,lem:exists_pow_mul_pow_le_exp_sq_mul_sqrt_log}
The rows of $B$ are the coefficients of the reduced polynomials of Lemma~\ref{lem:exists_iteratedDeriv_reduced_presentation}, bounded by Lemma~\ref{lem:exists_pow_mul_pow_le_exp_sq_mul_sqrt_log}; $M$ is taken large enough for twice as many unknowns as equations.
\end{proof}

% Prove2Me: FourExp.siegel_aux -- https://prove2.me/theorems/b421c5b8-a7f6-4c18-a1c0-d45a9e8fb2c7
\begin{lemma}[Solving the system]
\label{lem:siegel_aux}
\lean{Diaz.siegel_aux}
\leanok
Let $\omega,\omega_1\in\C$ and $Q\in\Z[X][Y]$ with $d=\deg_YQ\ge1$, such that no non-zero $A\in\Z[X][Y]$ with $\deg_YA<d$ vanishes at $(\omega,\omega_1)$, and let $\kappa_1>0$. There is $\kappa\ge\kappa_1$ such that for every large $N$, every $M$ with $0<M\le\kappa_1S$, and every $R\times U$ integer matrix $B$ with $U=S(2N)^{2}Md$, $2R\le U$ and entries of absolute value at most $e^{\kappa_1N^{2}\sqrt{\log N}}$, there is an integer vector $q$ with $Bq=0$ such that every $|q(i,j,h,\mu,\nu)|\le e^{\kappa N^{2}\sqrt{\log N}}$, some $c_q(i,j,h)\neq0$, and every $|c_q(i,j,h)|\le e^{\kappa N^{2}\sqrt{\log N}}$.

{\small\emph{Source:} Siegel's lemma \cite[Ch.~I, \S2]{Lang66}, as used in \cite[Lemme 4]{W73}.}
\end{lemma}

\begin{proof}
\leanok
\uses{lem:siegel_entrywise}
Lemma~\ref{lem:siegel_entrywise} gives $q\neq0$. A non-zero block $q(i,j,h,\cdot,\cdot)$ is a polynomial of $Y$-degree less than $d$, so by minimality $c_q(i,j,h)\neq0$.
\end{proof}

% Prove2Me: FourExp.auxiliary_function_alg -- https://prove2.me/theorems/88de0071-cb97-40ae-90e2-33e7eb593322
\begin{lemma}[The auxiliary function]
\label{lem:auxiliary_function_alg}
\lean{Diaz.auxiliary_function_alg}
\leanok
Let $x_1,x_2,y_1,y_2\in\C$ have presentation data, and suppose every $e^{x_iy_j}$ is algebraic. There is $\kappa>0$ such that for every large $N$ there are $M\le\kappa S$ and integers $q(i,j,h,\mu,\nu)$ ($i<S$, $j,h<2N$, $\mu<M$, $\nu<d$) of absolute value at most $e^{\kappa N^{2}\sqrt{\log N}}$, such that the coefficients $c_q(i,j,h)$ are not all zero and of absolute value at most $e^{\kappa N^{2}\sqrt{\log N}}$, and
\[ F_q^{(m)}(ay_1+by_2)=0\qquad\text{for all }a<t_1,\ b<t_2,\ m<S. \]

{\small\emph{Source:} \cite[Lemme 4]{W73}, with Siegel's lemma.}
\end{lemma}

\begin{proof}
\leanok
\uses{lem:aux_linear_system,lem:siegel_aux}
Lemma~\ref{lem:aux_linear_system}, then Lemma~\ref{lem:siegel_aux}.
\end{proof}

% Prove2Me: FourExp.extrapolation_numbers -- https://prove2.me/theorems/4613db26-460f-498c-a106-b0b8709cc575
\begin{lemma}[The numbers of the extrapolation]
\label{lem:extrapolation_numbers}
\lean{Diaz.extrapolation_numbers}
\leanok
Let $X\ge0$, $Y\ge1$ and $\kappa\ge0$ be real, and let $N,S,t_1,t_2,s$ be natural numbers with $s\le S$ and $N\ge32(13+\kappa+\log Y+2XY)$. Write $w=\sqrt{\log N}$ and suppose $Sw\le N^{2}\le2Sw$, $N\le2t_1w$ and $Nw\le2t_2$. Then
\[ 2\,s!\,(N-1)^{-t_1t_2S}\cdot S(2N)^{2}e^{\kappa N^{2}w}\,(YN^{2}w)^{S}\,e^{2NX\cdot YN^{2}w}\le e^{-N^{4}w/32}. \]

{\small\emph{Source:} The parameter estimate in the proof of \cite[Lemme 5]{W73}.}
\end{lemma}

\begin{proof}
\leanok
\end{proof}

% Prove2Me: FourExp.extrapolation -- https://prove2.me/theorems/478273e3-7383-4cf2-98e8-dd6ab2db30f3
\begin{lemma}[Extrapolation]
\label{lem:extrapolation}
\lean{Diaz.extrapolation}
\leanok
Let $x_1,x_2\in\C$, let $y_1,y_2$ be linearly independent over $\Q$, and let $\kappa>0$. There are $\kappa'>0$ and $N_0$ such that for every $N>N_0$ and all coefficients with $|c(i,j,h)|\le e^{\kappa N^{2}\sqrt{\log N}}$: if $F^{(m)}(ay_1+by_2)=0$ for all $a<t_1$, $b<t_2$ and $m<S$, then
\[ |F^{(s)}(ay_1+by_2)|\le\exp\bigl(-N^{4}\sqrt{\log N}/\kappa'\bigr)\qquad\text{for all }s<S',\ a<R_1,\ b<R_2. \]

{\small\emph{Source:} \cite[Lemme 5]{W73}.}
\end{lemma}

\begin{proof}
\leanok
\uses{lem:expPoly_grid_estimate,lem:extrapolation_numbers}
Lemma~\ref{lem:expPoly_grid_estimate} for the grid $\{ay_1+by_2:a<t_1,\ b<t_2\}$, with the numbers of Lemma~\ref{lem:extrapolation_numbers}.
\end{proof}

% Prove2Me: FourExp.exists_int_norm -- https://prove2.me/theorems/3b96e024-7747-4dc5-92bd-7e2392eba335
\begin{lemma}[Norms down to $\Q(\omega)$]
\label{lem:exists_int_norm}
\lean{Diaz.exists_int_norm}
\leanok
Let $\omega,\omega_1\in\C$ and let $Q\in\Z[X][Y]$ be monic in $Y$ of degree $d$, with $Q(\omega,\omega_1)=0$ and minimal there. There is $c\in\N$ such that for every $A\in\Z[X][Y]$ of $Y$-degree less than $d$, length at most $b$ and $X$-degree at most $e$, with $A(\omega,\omega_1)\neq0$, there is a non-zero $P\in\Z[X]$ of length at most $(cb)^{d}$ and degree at most $d(e+c)$ with
\[ |P(\omega)|\le|A(\omega,\omega_1)|\,\bigl(cb\max(1,|\omega|)^{e}\bigr)^{d}. \]

{\small\emph{Source:} The core of the proof of \cite[Lemme 7]{W73}.}
\end{lemma}

\begin{proof}
\leanok
$P$ is the determinant of multiplication by $A$ on $\Z[X][Y]/(Q)$ in the basis $1,Y,\dots,Y^{d-1}$: the norm from $\Q(\omega,\omega_1)$ to $\Q(\omega)$, taken without building the field extension.
\end{proof}

% Prove2Me: FourExp.norm_to_polynomial_alg -- https://prove2.me/theorems/7f85aec3-0e4c-43b2-9ddb-b8c54d207c1d
\begin{lemma}[From a small value to a small polynomial]
\label{lem:norm_to_polynomial_alg}
\lean{Diaz.norm_to_polynomial_alg}
\leanok
Let $x_1,x_2,y_1,y_2\in\C$ have presentation data, suppose every $e^{x_iy_j}$ is algebraic, and let $\kappa,\kappa'>0$. There is $k>0$ such that for every $C$ there is $N_0$ with the following property. Let $N>N_0$, $M\le\kappa S$, and let $q(i,j,h,\mu,\nu)$ be integers of absolute value at most $e^{\kappa N^{2}\sqrt{\log N}}$. If $a<R_1$, $b<R_2$, $s<S'$ and
\[ 0<|F_q^{(s)}(ay_1+by_2)|\le\exp\bigl(-N^{4}\sqrt{\log N}/\kappa'\bigr), \]
then some $P\in\Z[X]$ is small at $\omega$ for $(N,C)$ with respect to the growth functions of constant $k$.

{\small\emph{Source:} \cite[Lemme 7]{W73}.}
\end{lemma}

\begin{proof}
\leanok
\uses{lem:exists_iteratedDeriv_reduced_presentation,lem:exists_int_norm,lem:exists_pow_mul_pow_le_exp_sq_mul_sqrt_log,lem:length_sum_le,lem:length_mul_le}
Lemma~\ref{lem:exists_iteratedDeriv_reduced_presentation} writes the value, times a power of $D(\omega,\omega_1)$, as $A(\omega,\omega_1)$; take for $P$ the norm of Lemma~\ref{lem:exists_int_norm}. The sizes come from Lemma~\ref{lem:exists_pow_mul_pow_le_exp_sq_mul_sqrt_log}.
\end{proof}

% Prove2Me: FourExp.construction_core_1973 -- https://prove2.me/theorems/41291ef0-fcbe-4551-92c0-bddd04ee2f84
\begin{lemma}[The core of the construction]
\label{lem:construction_core_1973}
\lean{Diaz.construction_core_1973}
\leanok
Let $x_1,x_2$ be linearly independent over $\Q$, and likewise $y_1,y_2$. Suppose every $e^{x_iy_j}$ is algebraic and $\trdeg_{\Q}\Q[x_1,x_2,y_1,y_2]\le1$. Then there are a transcendental $\omega\in\C$ and $k>0$ such that for every $C$ and every large $N$ there are coefficients $c(i,j,h)$, not all zero, with the following property: whenever $a<R_1$, $b<R_2$, $s<S'$ and $F^{(s)}(ay_1+by_2)\neq0$, some $P\in\Z[X]$ is small at $\omega$ for $(N,C)$ with respect to the growth functions of constant $k$.

{\small\emph{Source:} \cite[\S III, (4) and Lemmes 4, 5, 7]{W73}.}
\end{lemma}

\begin{proof}
\leanok
\uses{lem:trdeg_one_presentation,lem:auxiliary_function_alg,lem:extrapolation,lem:norm_to_polynomial_alg}
Lemma~\ref{lem:trdeg_one_presentation} gives presentation data and Lemma~\ref{lem:auxiliary_function_alg} the auxiliary function. Lemma~\ref{lem:extrapolation} makes its derivatives small on the larger grid, and Lemma~\ref{lem:norm_to_polynomial_alg} turns a non-zero value into $P$.
\end{proof}

% Prove2Me: FourExp.construction_growth -- https://prove2.me/theorems/756ef0b1-6367-48de-9b97-b031f9277369
\begin{lemma}[The growth functions]
\label{lem:construction_growth}
\lean{Diaz.construction_growth}
\leanok
For every $k>0$ the growth functions $\sigma_1,\sigma_2$ are strictly increasing and tend to $+\infty$, and for every $x>0$
\[ \sigma_2(x)\le\sigma_1(x),\qquad\sigma_i(x+1)\le3\,\sigma_i(x)\quad(i=1,2). \]

{\small\emph{Source:} Elementary; the functions of the end of \cite[\S III]{W73}.}
\end{lemma}

\begin{proof}
\leanok
\end{proof}

% Prove2Me: FourExp.construction_count_1973 -- https://prove2.me/theorems/fede9f8c-c982-4629-9b09-f945e6eb22f6
\begin{lemma}[The parameters beat the zero count]
\label{lem:construction_count_1973}
\lean{Diaz.construction_count_1973}
\leanok
Let $X,Y_1,Y_2\ge0$. There is $N_0$ such that for every integer $N>N_0$, with $T=2N$, $\lambda=1/20$ and $n=ST^{2}$,
\[ \frac{n}{\lambda}+2\,\frac{1+n^{\lambda}}{\lambda\log n}\bigl(1+(R_1Y_1+R_2Y_2)\,TX\bigr)\le R_1R_2S'. \]

{\small\emph{Source:} The parameters of \cite[(4) and Lemme 6]{W73} against the zero estimate \cite[\S4, Lemme 3]{W71}.}
\end{lemma}

\begin{proof}
\leanok
$n/\lambda$ is about $80N^{4}/\sqrt{\log N}$ and $R_1R_2S'$ about $98N^{4}/\sqrt{\log N}$, while the second term is $O(N^{2.2+\varepsilon})$.
\end{proof}

% Prove2Me: FourExp.auxiliary_construction -- https://prove2.me/theorems/98a064ef-3a62-440c-94b0-56817780aca6
\begin{lemma}[The construction]
\label{lem:auxiliary_construction}
\lean{Diaz.auxiliary_construction}
\leanok
Let $l_{11},l_{12},l_{21},l_{22}$ satisfy the hypotheses of Lemma~\ref{lem:rank_one_parametrization}. Then there are a transcendental $\omega$, functions $\sigma_1,\sigma_2$ and constants $a_1,a_2\ge1$ satisfying the hypotheses of Theorem~\ref{thm:transcendence_criterion}, and pairs $x_1,x_2$ and $y_1,y_2$, each linearly independent over $\Q$, with the following property. For every $C$ and every large $N$ there are natural numbers $S,T,R_1,R_2,S'$, coefficients $c(i,j,h)\in\C$ ($i<S$, $j,h<T$), not all zero, and $\lambda>0$ such that
\begin{itemize}
\item the inequality of Lemma~\ref{lem:nonvanishing_derivative} holds, and
\item whenever $a<R_1$, $b<R_2$, $s<S'$ and $G^{(s)}(ay_1+by_2)\neq0$, where
\[ G(z)=\sum_{i<S}\sum_{j,h<T}c(i,j,h)\,z^{i}e^{(jx_1+hx_2)z}, \]
some $P\in\Z[X]$ is small at $\omega$ for $(N,C)$ with respect to $\sigma_1,\sigma_2$.
\end{itemize}

{\small\emph{Source:} \cite[\S III, (4) and Lemmes 4, 5, 7]{W73}.}
\end{lemma}

\begin{proof}
\leanok
\uses{lem:rank_one_parametrization,lem:construction_growth,lem:construction_count_1973,lem:construction_core_1973}
Lemma~\ref{lem:rank_one_parametrization} gives $x$ and $y$, and Lemma~\ref{lem:construction_core_1973} gives $\omega$ and the coefficients, with the parameters of the notation above, the growth functions of Lemma~\ref{lem:construction_growth} ($a_1=a_2=3$) and the count of Lemma~\ref{lem:construction_count_1973}.
\end{proof}

% Prove2Me: FourExp.small_polynomials_of_counterexample -- https://prove2.me/theorems/81f3c358-88a8-46e9-a5af-5c2c359b63f6
\begin{lemma}[A counterexample gives small polynomials]
\label{lem:small_polynomials_of_counterexample}
\lean{Diaz.small_polynomials_of_counterexample}
\leanok
Let $l_{11},l_{12},l_{21},l_{22}$ be non-zero complex numbers with every $e^{l_{ij}}$ algebraic, $l_{11}l_{22}=l_{12}l_{21}$ and $\trdeg_{\Q}\Q[l_{11},l_{12},l_{21},l_{22}]\le1$, and suppose that neither the rows nor the columns of $(l_{ij})$ are linearly dependent over $\Q$. Then there are a transcendental $\omega\in\C$, functions $\sigma_1,\sigma_2$ and constants $a_1,a_2\ge1$ satisfying the hypotheses of Theorem~\ref{thm:transcendence_criterion}, such that for every $C\in\R$ there is $N_0$ and for every $N>N_0$ a polynomial $P_N\in\Z[X]$ that is small at $\omega$ for $(N,C)$ with respect to $\sigma_1,\sigma_2$.

{\small\emph{Source:} \cite[\S III, Lemmes 4--7]{W73}, with the zero estimate of \cite[\S4, Lemme 3]{W71}; see also \cite{Bro74}.}
\end{lemma}

\begin{proof}
\leanok
\uses{lem:auxiliary_construction,lem:nonvanishing_derivative}
Lemma~\ref{lem:auxiliary_construction}, with a non-zero derivative supplied by Lemma~\ref{lem:nonvanishing_derivative}.
\end{proof}

\section{The theorem}

% Prove2Me: DiazModulus.four_exponentials_trdeg_one -- https://prove2.me/theorems/2c0f35ea-b824-4ce3-a701-01b48dc27a97
\begin{theorem}[Four exponentials in transcendence degree one]
\label{thm:four_exponentials_trdeg_one}
\lean{Diaz.four_exponentials_trdeg_one}
\leanok
Let $l_{11},l_{12},l_{21},l_{22}$ be non-zero logarithms of algebraic numbers with
\[ l_{11}l_{22}=l_{12}l_{21},\qquad\trdeg_{\Q}\Q(l_{11},l_{12},l_{21},l_{22})\le1. \]
Then the two rows, or the two columns, of $\begin{pmatrix}l_{11}&l_{12}\\ l_{21}&l_{22}\end{pmatrix}$ are linearly dependent over $\Q$. Equivalently, if $x_1,x_2$ and $y_1,y_2$ are pairs of complex numbers linearly independent over $\Q$ and the four products $x_iy_j$ are logarithms of algebraic numbers, then $\trdeg_{\Q}\Q(x_1,x_2,y_1,y_2)\ge2$.

{\small\emph{Source:} Waldschmidt \cite[Cor.~4]{W73} and Brownawell \cite[Cor.~7]{Bro74}; stated in this form as \cite[Theorem 1]{RW95}.}
\end{theorem}

\begin{proof}
\leanok
\uses{thm:transcendence_criterion,lem:small_polynomials_of_counterexample}
Otherwise Lemma~\ref{lem:small_polynomials_of_counterexample} gives polynomials too small at a transcendental $\omega$, and Theorem~\ref{thm:transcendence_criterion} makes $\omega$ algebraic.
\end{proof}
